A further tension in the literature concerns what ``accuracy'' means when papers are compared. Elicitation-based BN studies often evaluate risk classification on scarce projects or on resampled expert cases~\cite{Khodabakhshian2025,Comu2023}. Progress-prediction studies evaluate stage or milestone labels with multinomial models~\cite{Tagharobi2025}. Delay-cause surveys report prevalence rankings without predictive hold-outs~\cite{Assaf2006,Sambasivan2007,Doloi2012}. Treating their headline accuracies as interchangeable leaderboard entries would be methodologically careless. The appropriate use of those numbers in the present paper is as published reference points on the accuracy scale that practitioners already cite, accompanied by an explicit scope caveat and by additional metrics (ROC/PR, Brier, ECE, net benefit) that many of those papers do not report~\cite{Egwim2021,Gneiting2007,Vickers2006,Saito2015}.

VoI research in reliability engineering likewise cannot be copied verbatim into construction delay without redefining actions and observations. Component inspection typically reveals degradation measurements that enter a physics-informed likelihood~\cite{Yuan2021,Fauriat2020}. Sensor placement VoI optimizes spatial information under Gaussian or related field models~\cite{Malings2016}. Binary-system VoI analyses assume structurally defined cut sets~\cite{Song2021}. Construction delay-factor inspection, by contrast, reveals administrative presence/absence of a hypothesized cause, and the downstream action is often an alert or escalation rather than an immediate repair. The mathematics of expected loss reduction still applies~\cite{Howard1966,Memarzadeh2016}, but the empirical behavior of VoI policies depends on how sparse and how decisive those administrative observations are. That dependence is exactly what the NYC budgeted scoreboard measures.

Calibration-centric forecasting theory supplies the last connecting tissue. A model can achieve high AUC while producing probabilities that are too sharp or too timid for thresholded decisions~\cite{Gneiting2007,Murphy1973}. Strictly proper scoring rules penalize those failures~\cite{GneitingRaftery2007,Brier1950}. Decision curves translate probability quality into net benefit as a function of threshold preference~\cite{Vickers2006}. Construction AI reviews celebrate algorithmic diversity~\cite{Darko2020,PanYZhang2021,Abioye2021,Regona2024} but comparatively rarely insist on this probabilistic evaluation stack. Positioning DRIN against that omission is as important as positioning it against BN delay classics: the contribution is joint methodological coupling plus evaluation discipline, not a claim that boosting or Noisy-OR are newly invented.
